\section{What adaptive evaluation costs}
\label{sec:cost}

Adaptive testing saves items only after a calibration panel has answered the whole bank.
Counting one administered item as the unit of cost, evaluating $M$ new systems costs $MI$
with the full bank of $I$ items and $NI+MT$ with a fresh panel of $N$ calibrators and tests of
$T$ items, so adaptive testing pays off once $M>NI/(I-T)$. Table~\ref{tab:cost} gives, for
$N=8$ and $T=100$, this break-even point (between \CostBreakFullMin\ and \CostBreakFullMax\
new systems), the number of items with which EOST-M matches the score error of 100 items
sampled uniformly and scored by the sample mean (Naive) or by the same calibrated model
(Random:mg), and the break-even point against Naive. Uniform subsampling needs no panel, so
it is the fairer comparison: EOST-M saves about \CostSaveNaive\ items per system against it
and breaks even after a median of \CostBreakNaiveMed\ new systems (from \CostBreakNaiveMin\
to \CostBreakNaiveMax). When the calibrators' results already exist---public leaderboards or
earlier internal evaluations---the panel is free and adaptive testing is cheaper from the
first system; when items are expensive, as in agentic benchmarks, the per-system saving
dominates quickly. Calibrating a panel and running a 100-item test take seconds on one CPU
core.

\begin{table}[htbp]
\centering\footnotesize
\caption{Cost of adaptive evaluation with a fresh panel of eight calibrators.}
\label{tab:cost}
\begin{tabular}{@{}lrrrrr@{}}
\toprule
 & & break-even vs & items to match & break-even vs & items to match \\
Benchmark & $I$ & full evaluation & Naive@100 & Naive & Random:mg@100 \\
\midrule
GSM8K & 1319 & 8.7 & 70 & 352 & 80 \\
ARC-C & 1172 & 8.7 & 20 & 117 & 90 \\
HellaSwag & 2000 & 8.4 & 8 & 174 & 25 \\
WinoGrande & 1267 & 8.7 & 20 & 127 & 60 \\
MMLU & 2000 & 8.4 & 15 & 188 & 100 \\
SWE-bench V. & 500 & 10.0 & 40 & 67 & 70 \\
LiveCodeBench & 1055 & 8.8 & 25 & 113 & 60 \\
CRUXEval-O & 800 & 9.1 & 25 & 85 & 60 \\
\bottomrule
\end{tabular}
\end{table}


\section{Discussion}
\label{sec:discussion}

\paragraph{Small-panel calibration uncertainty is the problem.} With four or eight
calibrators a large fraction of items is answered identically by every calibrator, so their
difficulty is located only by the prior. Plug-in testing treats these
poorly located parameters as known: Fisher selection credits such items with information
they do not carry, the plug-in posterior becomes overconfident, and the predicted outcomes of
the unobserved items inherit the parameter errors---the bias floor of
Figure~\ref{fig:curves-err}. The expected response function caps the information of such
items (Proposition~\ref{prop:cap}) and predicts their outcomes with the uncertainty they
actually have. Both effects fade as the panel grows (Figure~\ref{fig:heatmap}): for the very
large panels of recent work \citep{li2026atlas}, plug-in Fisher information is an adequate
approximation; for new, private or expensive benchmarks, self-generated item pools, or the
panels of four to eight systems that a single laboratory can afford, it is not.

\paragraph{Integrating and constraining are two remedies.} Integrating the item posterior
(EOST-M) predicted better than constraining the discriminations to a common value (C1), more
so with sixteen calibrators, whereas the constrained model's intervals were better calibrated
because they are wider, not because its estimates are more accurate. A small panel must
therefore not be trusted with item-specific discriminations: integrating their uncertainty is
preferable when prediction and score accuracy matter most, and constraining them gives a
simpler and more conservative test.

\paragraph{Why not always use a fully Bayesian adaptive test?} It is the principled answer
to calibration uncertainty, and \citet{fink2025accounting} found it the most accurate
correction with hundreds of calibration responses per item. With eight calibrators it did not
help measurably, although equivalence was not shown either. By Lemma~\ref{lem:exact}, EOST-M
is the fully Bayesian model with the calibrators' abilities and the hyperprior held fixed; in
this regime what matters, as far as our data can tell, is the item-level uncertainty that
both integrate. A fully Bayesian test also needs a Markov chain per panel, whose tuning is a
source of error in its own right: our frozen burn-in was too short on some banks, and even an
exact sampler with collapsed moves did not meet our convergence criterion. EOST-M is the
simpler default. Two leads remain open: the near-converged sampler suggests a possible
advantage of Bayesian integration for estimating the score, and panel-level uncertainty may
matter more when the calibrators' abilities are poorly determined, as with short banks.

\paragraph{The criterion matters less than first hypothesized.} The pilot program and
study~1 agree on the direction of the original thesis---EIG learns more about competence than
predictive entropy---but bound its size: under the same calibration-marginal model the gain
on held-out prediction is small, and the information criteria are practically
interchangeable (a descriptive comparison; the failure of H1 is not a test of equivalence).
Corollary~\ref{cor:blind} explains why: entropy and information disagree only through the
spread of discriminations among items with similar success probabilities, which the expected
response function shrinks, and with a common discrimination the difference vanishes. For the
bank score, Proposition~\ref{prop:evr} shows that uncertainty sampling becomes the right
one-step criterion once competence is pinned down relative to the unobserved bank, and the
score-error curves show both regimes. What an adaptive test for LLMs must get right is the
treatment of calibration uncertainty, not the choice among information criteria.

\paragraph{Honest uncertainty.} None of the four remedies tried in study~2---integrating
the calibrators' abilities, leave-one-calibrator-out recalibration, mixing random items into
the test, recalibrating a constrained model---gave nominal coverage with intervals as narrow as
EOST-M's. The coverage decomposition and the prior predictive checks locate the problem where
adaptive selection meets a model that is misspecified in the dependence among real systems'
verdicts. Recalibration is not yet a validated remedy: it raised coverage, but neither frozen
criterion for it was met, and on the well-specified synthetic bank it moved coverage away
from the nominal level. Practitioners should report model-based intervals with the panel size
and their expected under-coverage, and use a Wilson interval from uniformly sampled items when
nominal coverage matters more than width.

\paragraph{Adaptive testing or benchmark reduction?} With a handful of calibrators,
tinyBenchmarks, ATLAS-style CAT and MINCE rely on quantities that a small panel cannot
determine well. In this regime our evidence favours adaptive selection under calibration
uncertainty, in the same direction on every benchmark but conclusively by the frozen
criterion only against MINCE; with large panels the differences should shrink.

\paragraph{Limitations.} (i) All evidence replays recorded verdicts: replay is exact for the
recorded data but omits the sampling variability of stochastic decoding, and the protocol we
release for live evaluation could not be run in our environment; the temporal validation
approximates, but does not replace, a prospective deployment. (ii) LiveCodeBench and
CRUXEval-O carry binarized pass rates; graded responses would need another measurement model.
(iii) The unidimensional 2PL is misspecified in the dependence among real systems' verdicts.
(iv) Strict clusters are a heuristic that may miss distant relatives such as distilled models,
although R2 and C1 did not depend on how they were built. (v) Recalibration rests on very
few pseudo-targets, without a finite-sample guarantee. (vi) A fully converged Bayesian
comparator, for instance by Hamiltonian Monte Carlo, remains to be run. (vii) The analysis
plans were frozen by hashing and committing them, without an external time stamp. (viii) Item
exposure was not controlled; an adaptive benchmark that serves the same informative items to
every model would accelerate their contamination \citep{wainer2000cat}. (ix) Joint MAP
calibration made the planned one-parameter sensitivity analysis degenerate; marginal maximum
likelihood would avoid this. (x) Self-generated items were studied only in a historical
pilot. (xi) The cross-benchmark meta-analysis is secondary and limited by shared systems.

\paragraph{Recommendations.} With a small calibration panel: (1) do not plug in
item-specific discriminations calibrated on fewer than a few dozen systems, but integrate their
uncertainty with an empirical-Bayes prior, or constrain them when a simple, conservative test
matters more than accuracy; (2) select by any information criterion computed from the marginal
model, and let the score-targeted criterion take over late in the test when the estimand is
the score; (3) prefer the calibration-marginal model to the Bayesian treatments of the
calibrators that we ran, which made no significant difference; (4) report model-based
intervals with the panel size and their expected under-coverage; (5) exclude the target's whole
cluster from its panel and analyse with clusters as units; (6) count the panel in the cost.

\section{Conclusion}
\label{sec:conclusion}

Adaptive testing can make LLM evaluation cheaper, but with the small calibration panels of
new, private or expensive benchmarks it works only if the uncertainty of the item calibration
is taken seriously. Integrating it through expected response functions with an empirical-Bayes
prior is simple, exact given the calibrators, and caps the information of items that a small
panel cannot locate. Under frozen analysis plans, this predicted held-out verdicts better than
plug-in Fisher-information testing and a regularized Rasch-type test, estimated scores more
accurately than uniform subsampling and MINCE, was ahead of random items, ATLAS-style testing
and tinyBenchmarks on every benchmark without meeting the frozen criterion for these
comparisons, and was not significantly different from Bayesian CATs that also integrate the
calibrators' abilities, without being shown equivalent. The thesis that information rather than entropy should be maximized
holds in sign but is modest once uncertainty is integrated. The open problem is honest
uncertainty: model-based intervals remain too narrow, and recalibration from the panel improved
coverage only at the price of width.
